% Generated draft: requires mathematical and source review.
\documentclass[11pt]{article}
\usepackage[margin=25mm]{geometry}
\usepackage{fontspec}
\setmainfont{Latin Modern Roman}
\usepackage{amsmath,amssymb,mathtools,unicode-math}
\setmathfont{Latin Modern Math}
\usepackage{enumitem,needspace,xurl,hyperref}
\setlength{\parindent}{0pt}
\setlength{\parskip}{5pt}
\setlength{\emergencystretch}{3em}
\newcommand{\N}{\mathbb{N}}
\newcommand{\Z}{\mathbb{Z}}
\newcommand{\Q}{\mathbb{Q}}
\newcommand{\R}{\mathbb{R}}
\newcommand{\C}{\mathbb{C}}
\newcommand{\F}{\mathbb{F}}
\newcommand{\A}{\mathbb{A}}
\newcommand{\PP}{\mathbb{P}}
\newcommand{\E}{\mathbb{E}}
\newcommand{\eps}{\varepsilon}
\newcommand{\dd}{\,\mathrm{d}}
\begin{document}
\section*{UnsolvedMath ID: 30006323}
\textbf{Source problem code:} catalogue entry
\subsection{Definitions and mathematical statement}
The supplied record concerns the Plancherel algebra associated with a hyperspherical variety and its proposed dual Hamiltonian variety \(M^\vee\). The precise class of varieties, the base field, coefficient theory, and categorical framework are not specified in the supplied excerpt. The following formulation is therefore conditional on the intended framework and its existence.

Write \(X\) for the geometric object supporting the basic sheaf \(\delta_X\), \(\mathcal C_X\) for its intended sheaf category, and \(\mathcal H\) for the convolution category acting on it. With \(\star\) denoting this action, define
\[
\operatorname{act}_{\delta_X}\colon\mathcal H\longrightarrow\mathcal C_X,
\qquad F\longmapsto F\star\delta_X.
\]
Assuming that this functor has the right adjoint appearing in the source, define the internal endomorphism object
\[
A_X:=\operatorname{act}_{\delta_X}^{R}(\delta_X).
\]
The source calls this object the \emph{Plancherel algebra}, or the \emph{relative Coulomb branch algebra}, and refers to [3, Definition 8.1.4]. Its algebra structure is understood to be the internal endomorphism algebra structure in the intended categorical framework.

Let \(\operatorname{Sat}_{\mathrm{der}}\) denote the intended derived geometric Satake functor, and let \(\mathcal O(M^\vee)\) denote the function algebra of the proposed dual Hamiltonian variety. The supplied conjecture, attributed in the excerpt to [3, Conjecture 8.1.8], asks whether, for every instance satisfying the hypotheses of that conjecture, there is an identification
\[
\operatorname{Sat}_{\mathrm{der}}(A_X)
\simeq \mathcal O(M^\vee)
\]
with the intended algebra structures. The supplied record does not specify those hypotheses, the target category, or the equivariant, graded, or derived conventions needed to interpret this identification precisely; no unrestricted quantification over hyperspherical varieties is asserted here.

The excerpt states that this identification follows from a stronger local conjecture and is more tractable because the relevant sheaf category is not yet really defined. Neither the local conjecture nor a construction of that category is supplied. Consequently, the displayed assertion is a conditional statement of the intended conjecture, requiring the missing framework and hypotheses before it becomes a fully specified mathematical proposition.
\subsection{Short English statement}
The conjecture asks whether derived geometric Satake sends the internal endomorphism algebra of the basic sheaf to the function algebra of the dual Hamiltonian variety. The supplied excerpt leaves the categorical framework and hypotheses unspecified and explicitly qualifies the definition of the sheaf category. Its open status is supplied dated curation, not an independently verified conclusion.
\subsection{Sources}
\begin{itemize}
\item Supplied local UnsolvedMath registry record, permanent ID \(30006323\), title \emph{Plancherel Algebras of Dual Hyperspherical Varieties}, exact source problem code \texttt{catalogue entry}, in the Oberwolfach Reports Open Problems collection. No individual proposer is supplied; the recorded proposed year is 2025.
\item \emph{Arbeitsgemeinschaft: Relative Langlands Duality} (2025), Oberwolfach Reports, pp.\ 811--882, DOI 10.4171/owr/2025/18. Supplied report reference and source URL: \url{https://doi.org/10.4171/owr/2025/18}.
\item \emph{Relative Langlands Duality}, supplied literature reference: \url{https://arxiv.org/abs/2409.04677}. The report excerpt attributes the Plancherel algebra terminology to [3, Definition 8.1.4] and the proposed identification to [3, Conjecture 8.1.8]; the supplied excerpt does not explicitly resolve the bibliographic label [3].
\item Literature assessment dated exactly \texttt{2026-08-22}: the supplied curation labels the problem open and identifies the general assertion as Conjecture 8.1.8. It states that examples and formal reductions do not prove the assertion in the general spherical setting. This is supplied dated curation, not independently verified literature assessment. The listed literature sources are the OWR report and \emph{Relative Langlands Duality} linked above.
\item The registry labels its statement \texttt{corrected\_verified} and reports that the official report was checked and its notation repaired. These are supplied source-verification claims; no independent source visit or verification is asserted in this draft.
\end{itemize}
\end{document}
